\umaappendix{Detalle de la implementación}
\label{ap:implementacion}

Este apéndice completa el capítulo~\ref{cap:implementacion} con los detalles de implementación de cada módulo.

\section{Módulo de ingesta en Go}
\label{sec:impl-go-detalle}

\paragraph{Cliente HTTP y orquestador.} Las dos fuentes usan el cliente \texttt{httpx}, que limita el ritmo por sitio (dos peticiones por segundo al BOE y una a la Junta, por cortesía con servidores públicos), corta cada intento a los 30~s, rechaza respuestas de más de 50~MiB y reintenta hasta tres veces los errores de red, 429 y 5xx con espera exponencial. Todos los trabajos, del planificador, de la API HTTP o de la interfaz por gRPC, pasan por un orquestador que solo admite uno a la vez (RF2). Antes de descargar pide las reglas de filtrado a la API, y si no puede obtenerlas el trabajo falla (RF3); después recorre cada fuente y día del rango (31 como máximo), saltando los días cuyo sumario falla, y procesa los documentos de un día con un máximo de cinco \emph{goroutines} a la vez.

\paragraph{De la descarga a la cola.} Cada documento se descarga, se filtra (en minúsculas y sin tildes, por organismo o por palabra clave en título o texto) y, si pasa, se guarda en S3 en dos objetos: el original, como copia de seguridad, y el JSON con el texto y los metadatos, que es obligatorio. Solo después de guardar el JSON se publica el mensaje con su ruta, así que un \emph{worker} nunca recibe un documento que no existe; si la publicación falla, el administrador puede encolarlo desde la interfaz (RF27). El publicador de RabbitMQ \cite{rabbitmq} espera la confirmación del \emph{broker} para cada mensaje y declara las colas con los mismos argumentos que .NET; cada mensaje lleva el contexto de la traza (sección~\ref{sec:impl-observabilidad}). En gRPC, el \emph{scraper} es cliente (pide las reglas) y servidor (recibe \texttt{ForceScrape}), siempre con la clave interna comparada en tiempo constante. El planificador lanza a la 1:00 la descarga del día anterior de las fuentes configuradas (\texttt{BOE,BOJA} en producción), y toda la configuración se valida al arrancar devolviendo todos los errores juntos.

\section{Ingesta en .NET}
\label{sec:impl-ingesta-detalle}

\paragraph{\emph{Workers} y consumidor base.} El proyecto \texttt{Worker} es un \emph{host} mínimo que registra el consumidor de su variante según \texttt{WORKER\_PIPELINE\_MODE}; en la ingesta no hay ayuntamiento de petición y los avisos se publican en RabbitMQ en lugar de SignalR (sección~\ref{sec:ingesta-diseno}). Toda la lógica común de RabbitMQ está en \texttt{RabbitMqConsumerBase<T>}: se reconecta, recibe como máximo cinco mensajes sin confirmar, abre un ámbito de inyección de dependencias por mensaje y aplica la política de fallos del diseño mediante una función pura, \texttt{Decide}, que se prueba sin \emph{broker}. Un JSON ilegible o una \texttt{NonRetryableMessageException} van directos a la DLQ; cualquier otra excepción se reintenta hasta tres veces a través de una cola \texttt{.retry} sin consumidor, cuyo TTL de 30~s devuelve el mensaje a la cola de trabajo.

\paragraph{Nuevas estrategias de fragmentación.} Las estrategias del estudio son dos clases más de \texttt{IChunkingStrategy}. \texttt{RecursiveChunkingStrategy} prueba los separadores en orden (párrafo, salto de línea, punto, punto y coma, dos puntos y espacio), vuelve a dividir con el siguiente lo que no cabe y une piezas consecutivas hasta 400 \emph{tokens}, con 50 de solapamiento formado por piezas completas. \texttt{SemanticChunkingStrategy} divide el texto en frases con un separador que respeta las abreviaturas de los boletines (\emph{art.}, \emph{núm.}, \emph{Sr.}\ldots), calcula el \emph{embedding} de cada frase con sus vecinas y corta donde la distancia coseno entre ventanas consecutivas supera el percentil 95 del documento, uniendo los fragmentos de menos de 100 \emph{tokens} y dividiendo con la recursiva los de más de 400. Las dos reutilizan entera la \emph{pipeline} básica, registrada una vez por variante con su estrategia y su colección, de modo que la ingesta, la búsqueda y la generación son el mismo código y solo cambian los fragmentos (sección~\ref{sec:pipelines}).

\section{Orquestación RAG}
\label{sec:impl-rag-detalle}

\paragraph{Del \emph{hub} a la pipeline.} El método \texttt{StreamMessage} del \emph{hub} de chat devuelve un \texttt{IAsyncEnumerable<string>}, que SignalR \cite{signalr} transmite al cliente como un flujo: cada \texttt{yield return} de la pipeline llega al navegador como una parte de la respuesta. El \emph{hub} construye el turno (\texttt{ChatTurn}) con la pregunta, el historial de la sesión, el usuario y su ayuntamiento, y se lo pasa a \texttt{AiChatService}, que elige la pipeline como se explicó en la sección~\ref{sec:pipelines}. Si el cliente cancela el flujo (cierra la pestaña o pulsa «detener»), el \texttt{CancellationToken} llega hasta la llamada al modelo y la corta. Cuando la respuesta termina, el \emph{hub} decide si lanzar el \emph{shadow testing}, que se ejecuta en la cola de trabajo en segundo plano con el ayuntamiento del turno. Esa cola, \texttt{BackgroundWorkQueue}, es un canal acotado a 200 tareas que procesa cada una en su propio ámbito de inyección de dependencias; si está llena, la prueba se descarta con un aviso en el registro en lugar de bloquear el chat.

La fusión de las dos búsquedas de la \emph{pipeline} jerárquica, con $k=60$, es la del listado~\ref{lst:rrf}.

\begin{lstlisting}[language={[Sharp]C},caption={Fusión RRF de las dos búsquedas (\texttt{HybridRetrievalService})},label={lst:rrf}]
internal const int RrfK = 60;

internal static IReadOnlyList<(long Id, double Score)> Fuse(
    IReadOnlyList<long> denseRanking, IReadOnlyList<long> sparseRanking, int limit)
{
    var scores = new Dictionary<long, double>();
    foreach (var ranking in new[] { denseRanking, sparseRanking })
        for (var i = 0; i < ranking.Count; i++)
            scores[ranking[i]] = scores.GetValueOrDefault(ranking[i]) + 1.0 / (RrfK + i + 1);

    return scores.OrderByDescending(kvp => kvp.Value).Take(limit)
                 .Select(kvp => (kvp.Key, kvp.Value)).ToList();
}
\end{lstlisting}

Con los identificadores ganadores, una sola consulta de EF Core carga los fragmentos con su documento padre, y \texttt{ClearLanguageGenerationService} monta el \emph{prompt} con los padres en orden de relevancia dentro del presupuesto de caracteres descrito en la sección~\ref{sec:pipelines}.

\paragraph{\emph{Prompts}.} Todos los \emph{prompts} están en una única clase estática, \texttt{PromptCatalog}, y una prueba de arquitectura impide que los \emph{endpoints} y los \emph{hubs} definan los suyos (sección~\ref{sec:clean}). Tenerlos juntos facilita revisarlos y versionarlos con el código, y permite reproducir en el apéndice~\ref{ap:prompts} los de las \emph{pipelines} y la ingesta. Se escriben en español porque los documentos y las preguntas lo están, y cada uno indica el formato de salida que espera el código: texto libre para las respuestas, una pregunta por línea para el enriquecimiento y JSON para la expansión de consultas y el análisis del historial. Varias instrucciones se refieren a la pregunta como «del ciudadano» y piden tratarlo de tú, aunque quien usa el sistema es el personal municipal: las respuestas se redactan en lenguaje claro para que ese personal pueda trasladarlas tal cual a los vecinos, que son quienes suelen plantear estas dudas al ayuntamiento, y el enriquecimiento genera las preguntas que haría un ciudadano porque son las que luego llegan formuladas así.

\paragraph{Citas generadas por código.} Las primeras versiones pedían al modelo que terminara con la lista de fuentes, pero los modelos pequeños, como \texttt{llama3.2}, a menudo la omitían o inventaban los enlaces. Por eso, los \emph{prompts} actuales prohíben al modelo escribir esa sección, y es \texttt{SourceCitations} quien la añade al final del flujo, a partir de los documentos que realmente se enviaron al modelo (RF10). La sección solo se emite si la generación terminó sin errores ni cancelación. La lista indica qué documentos recibió el modelo, no cuáles respaldan cada frase: si el modelo descarta un documento porque no trata el tema de la pregunta, como le permite \texttt{StreamingUserTemplate} (apéndice~\ref{ap:prompts}), ese documento aparece igualmente en la lista. El enlace de los documentos del BOE apunta a su página en \texttt{boe.es}. El de los del BOJA se construye a partir de su identificador, que codifica la dirección de la que lo descargó el \emph{scraper}: \texttt{BOJA-2026-123-1} enlaza con \texttt{/boja/2026/123/1}. Si el identificador no sigue ese formato, el enlace apunta a la portada del boletín. En las ejecuciones del estudio todos los enlaces del BOJA llevaban todavía a la portada, pero eso no cambia las métricas, que no comprueban los enlaces.

\paragraph{Persistencia del chat.} Los mensajes no se escriben en la base de datos dentro de la petición, sino en un búfer en memoria (\texttt{ChatMessageBuffer}) que \texttt{ChatPersistenceWorker} vacía cada dos segundos. Si un lote falla, se reintenta mensaje a mensaje para que un mensaje problemático (por ejemplo, de una sesión borrada mientras tanto) no arrastre al resto. Los fallos transitorios, como una base de datos caída, mantienen los mensajes en el búfer sin gastar intentos, hasta 10.000 mensajes; los demás se reintentan tres veces y después se guardan en \texttt{FailedChatMessage}. Al parar el servicio se hace un último vaciado.

\section{Contexto personalizado y Arena}
\label{sec:impl-arena-detalle}

\paragraph{Preferencias del usuario.} Cada usuario tiene, en cada organización, una lista de temas y otra de zonas de interés (RF13, RF18), que se añaden al \emph{prompt} de sistema de la \emph{pipeline} básica con una instrucción que pide usarlas para priorizar o enfocar la respuesta. Se editan en el perfil o se sugieren a partir de las conversaciones (RF14): un análisis en segundo plano toma las cinco sesiones más recientes, las recorta a 8.000 caracteres y pide al modelo de la plataforma un JSON con temas y zonas, que se añaden sin borrar nunca lo que escribió el usuario.

\paragraph{Batallas y votos.} La página de la Arena ejecuta las dos \emph{pipelines} completas en paralelo y devuelve las dos respuestas; el modo Arena del chat reenvía las partes de cada una según llegan, etiquetadas como \emph{Alfa} o \emph{Beta}, leyendo de un canal por lado con \texttt{Task.WhenAny}. Cada \emph{pipeline} corre en su propio ámbito de inyección de dependencias, porque ni el \texttt{DbContext} ni los \emph{kernels} admiten uso concurrente. Qué \emph{pipeline} es \emph{Alfa} se decide al azar (RF21), y si un lado falla la batalla no se guarda. El voto (gana \emph{Alfa}, gana \emph{Beta}, empate o «las dos son malas») es definitivo (HTTP 409 si se repite); tras votar se revela qué \emph{pipeline} había detrás de cada opción (RF22), y cada batalla guarda el \texttt{TraceId} de su petición para enlazar el voto con sus trazas. Las batallas del \emph{shadow testing} se guardan igual, sin voto. \texttt{ArenaAnalyticsService} calcula por periodo las victorias, derrotas, empates, respuestas malas, tasa de victoria, latencia y \emph{tokens} de cada \emph{pipeline} (RF24); como recorre las batallas de todos los ayuntamientos, está en la lista revisada de la sección~\ref{sec:clean}, solo devuelve totales y es exclusivo del administrador de la plataforma (RF29).

\section{Interfaz en Angular}
\label{sec:impl-angular-detalle}

La interfaz es una aplicación Angular~21 \cite{angular} con componentes \emph{standalone}, organizada en \texttt{core} (autenticación, tiempo real, telemetría), \texttt{shared} y una carpeta por funcionalidad, cada una con sus servicios de acceso a la API. Todas las rutas se cargan de forma diferida, así que un usuario sin rol de administrador nunca descarga esas páginas; los guardas de ruta solo ocultan opciones y la autorización la decide la API (sección~\ref{sec:seguridad}). El estado del chat vive en \texttt{ChatStore}, con señales inmutables, de modo que una respuesta sigue llegando aunque el usuario cambie de sesión o de página (RF16); las partes del flujo se aplican una vez por fotograma (\texttt{requestAnimationFrame}) para no repintar con cada una. Las conexiones de SignalR obtienen el token de Auth0 de forma silenciosa y se reconectan solas, y los errores del flujo se traducen de los códigos estables de la API (sección~\ref{sec:proveedor-ia}). Las respuestas llegan en Markdown, se convierten con \texttt{marked} y se sanean con DOMPurify, que además abre los enlaces en una pestaña nueva sin acceso a la página de origen; las fuentes se muestran al final de cada mensaje, y en la Arena por separado para cada lado. La sesión se cierra tras 30 minutos sin actividad, pero no mientras se genera una respuesta (RF19), y en el móvil el menú lateral es un diálogo modal accesible (RF20).

\section{Orquestación y despliegue}
\label{sec:impl-despliegue-detalle}

En local, el \texttt{AppHost} de Aspire \cite{aspire} (sección~\ref{sec:arquitectura}) arranca todo el sistema con un solo comando y muestra registros, trazas y métricas sin más configuración. En producción, el sistema se ejecuta en servidores propios gestionados con Dokploy \cite{dokploy} (sección~\ref{sec:recursos}) a partir del \path{docker-compose.yml} del repositorio, que define RabbitMQ y Qdrant con volúmenes, la API, los \emph{workers} básico y jerárquico (las variantes del estudio solo arrancan con el perfil \texttt{chunking-study}), el \emph{scraper} y la pasarela, que es la única que publica un puerto; PostgreSQL se ejecuta en el propio servidor y el almacenamiento de objetos es Cloudflare R2. Las variables obligatorias usan la sintaxis \texttt{\$\{VAR:?mensaje\}}, que hace fallar el despliegue si falta alguna, y la API aplica las migraciones al arrancar.
% NOTA: producción usa Qwen3-Embedding-8B por OpenRouter (confirmado por Antonio el 19/11/2026); los valores por defecto del docker-compose.yml (google / text-embedding-004) se sustituyen con las variables AI_EMBEDDINGS_* de Dokploy.
Las imágenes .NET se generan con el soporte de contenedores del SDK (la de la pasarela incluye la compilación de Angular), y la del \emph{scraper}, con un \emph{Dockerfile} de dos etapas que acaba en una imagen \emph{distroless} sin privilegios. El flujo \texttt{deploy.yml} solo se lanza cuando la integración continua (sección~\ref{sec:calidad}) pasa en la rama principal: reconstruye solo los servicios que cambiaron, etiqueta todas las imágenes con \texttt{sha-<commit>} y llama al \emph{webhook} de Dokploy; volver atrás consiste en fijar \texttt{IMAGE\_TAG} a una versión anterior, teniendo en cuenta que las migraciones no se deshacen solas. Los pasos completos están en \texttt{docs/setup.md} y \texttt{docs/deploy-dokploy.md} del repositorio \cite{repositorio}.

\section{Observabilidad y coste}
\label{sec:impl-observabilidad-detalle}

Todos los servicios exportan trazas, métricas y registros con OpenTelemetry \cite{opentelemetry}: en .NET lo configura \texttt{ServiceDefaults}, y en Go el paquete \texttt{telemetry}, con registros JSON de \texttt{slog} que incluyen los identificadores de traza. El contexto de traza se propaga con \texttt{traceparent} por HTTP, gRPC y también por las cabeceras AMQP de RabbitMQ, de modo que una descarga lanzada desde la interfaz es una sola traza que va del navegador a la API, al \emph{scraper} y a los dos \emph{workers}, y los reintentos la conservan; una prueba de integración comprueba que .NET entiende la cabecera tal como la escribe Go.

Todas las llamadas a modelos pasan por uno de tres puntos (\texttt{ChatStreamingService}, la generación de la \emph{pipeline} jerárquica o el envoltorio \texttt{MeteredChatCompletion}), que registran la duración (con el código de error si falla), los \emph{tokens} de entrada y salida, el tiempo hasta la primera parte de la respuesta y el coste estimado, con los nombres de las convenciones de OpenTelemetry para IA generativa \cite{otelgenai} y el prefijo propio \texttt{app.} para el tiempo hasta el primer \emph{token} y el coste, que esas convenciones no definen. Cada medida lleva el proveedor, el modelo, el ayuntamiento, la \emph{pipeline}, la operación y si se usó el respaldo, nunca identificadores de usuario. El coste se calcula con una tabla de precios en euros por millón de \emph{tokens} por proveedor y modelo (\texttt{Ai:Pricing}); un modelo de pago sin precio registra cero y deja un aviso, y \texttt{TokenUsageExtractor} reúne la lectura de los \emph{tokens}, que cada conector devuelve de una forma, con pruebas por proveedor. La \emph{pipeline} jerárquica hace al menos dos llamadas por pregunta y envía un contexto mucho mayor que la básica, que hace una sola con tres fragmentos.
% TODO (Antonio): tabla de coste medio por pregunta para cada pipeline con modelos frontera (por ejemplo, Gemini u OpenAI vía OpenRouter) frente a modelos locales en Ollama: tokens de entrada y salida, coste en EUR y latencia p50, a partir del panel «LLM» de Grafana o de AiCallLog en un periodo concreto. Indicar el periodo y los precios usados (Ai:Pricing) con su fecha.

Los datos llegan a un OpenTelemetry Collector que reparte trazas, métricas y registros a Tempo, Mimir y Loki, visualizados en Grafana en un servidor propio del autor, al que los servicios llegan por un canal privado con un token. El panel «LLM» muestra el coste por ayuntamiento, modelo y operación, los percentiles de latencia, el uso del respaldo y los errores por código. En el navegador, Grafana Faro recoge errores y métricas de experiencia web y los envía a la pasarela, que los filtra y reenvía al Collector.
% VERIFICAR (Antonio): concretar el canal privado entre el VPS y el homelab (WireGuard, Tailscale u otro) y si ya está en uso en producción; docs/observability-architecture.md lo da como configuración recomendada.

\paragraph{Privacidad.} El contenido de las preguntas y respuestas no se envía a la telemetría (RNF9): \texttt{ServiceDefaults} activa solo el modo de diagnóstico de Semantic Kernel sin contenido, una prueba falla si el otro aparece en la configuración o en el código, y el Collector elimina esos atributos por si llegaran. Por eso se descartó Arize Phoenix \cite{arizephoenix} como herramienta específica de observabilidad de modelos (ADR~0002): funcionó sin cambios de código, pero sin el contenido de las conversaciones aportaba poco frente a Grafana. La prueba sí destapó un error: el extractor de \emph{tokens} no leía el formato de OpenAI, OpenRouter ni los \emph{tokens} de salida de Gemini, y el coste se registraba a cero; se corrigió con pruebas por proveedor.
